\section{Discussion and Outlook}
  \label{sec:discussion}

  \subsection{Scope of the Mechanism}
    \label{subsec:scope}

    The framework developed in this work
    proposes a structural interpretation of superconductivity
    as projective phase locking of $w=2$ composite configurations.
    It does not replace microscopic electronic Hamiltonians.
    Rather, it provides a geometric organizing principle
    that constrains admissible pairing channels,
    phase stiffness scaling,
    and symmetry selection
    across material families.

    In conventional superconductors,
    the mechanism reduces to mediator-stabilized composite formation,
    recovering the standard London and Ginzburg--Landau structure.
    In strongly correlated systems,
    pairing arises from the reduction of real-space frustration
    in fiber admissibility constraints,
    naturally separating amplitude and phase scales.

    The resulting picture unifies
    low-$T_c$ and high-$T_c$ regimes
    without introducing distinct pairing principles.
    Only the origin of the projective gap $\Delta_\Pi$
    differs between material classes.

  \subsection{Relation to Existing Theoretical Approaches}
    \label{subsec:relation-existing}

    The present framework differs from
    momentum-space interaction models
    in its real-space formulation.
    In particular,
    symmetry selection is determined
    by minimization of a fiber raccordement cost
    under lattice and frustration constraints,
    rather than by the structure of an interaction kernel
    $V(\mathbf{k},\mathbf{k}')$.

    This distinction does not invalidate
    spin-fluctuation or strong-coupling approaches.
    Instead,
    it constrains their admissible outcomes.
    If the experimentally observed gap symmetry
    is inconsistent with the frustration ratio
    in a given material,
    the geometric criterion proposed here
    would be falsified.

    The framework therefore operates
    at a structural level above specific microscopic models.
    Its role is to provide a symmetry and scaling constraint,
    not a detailed electronic mechanism.

  \subsection{Limitations}
    \label{subsec:limitations}

    Several limitations must be acknowledged.

    First,
    the cost functional $\mathcal{C}_{\Pi}$
    is not derived from a fully microscopic $\chi$-dynamics
    in the present work.
    Only its symmetry properties and leading scaling behavior
    are specified.
    A microscopic derivation remains an open program.

    Second,
    the proportionality between frustration density
    and magnetic structure factor
    is introduced as a first-order proxy.
    Although this assumption is physically motivated,
    non-linear corrections may alter numerical prefactors.
    The robustness of predictions therefore rests
    on qualitative rather than fine-tuned quantitative agreement.

    Third,
    the quasi-two-dimensional BKT description
    applies strictly in the anisotropic limit.
    Materials with strong interlayer coupling
    require a modified treatment,
    which may alter the detailed scaling of $T_c$.

  \subsection{Experimental Outlook}
    \label{subsec:experimental-outlook}

    The framework makes experimentally resolvable predictions.

    Most prominently,
    the predicted $s^\pm$ symmetry in nickelates
    and the intermediate sensitivity to non-magnetic disorder
    provide near-term tests.
    Phase-sensitive ARPES,
    impurity-substitution studies,
    and measurements of magnetic excitation spectra
    under pressure
    offer independent probes of the underlying mechanism.

    More broadly,
    systematic comparison of
    structure factor amplitudes,
    phase stiffness scaling,
    and gap symmetry
    across related material families
    can directly test the frustration-based selection principle.

  \subsection{Conceptual Implications}
    \label{subsec:conceptual-implications}

    If supported experimentally,
    the results suggest that superconductivity
    is not fundamentally defined
    by a specific pairing interaction,
    but by the existence of a stable composite class
    and the energetic cost of breaking global fiber coherence.

    In this view,
    pairing is a geometric necessity
    in regimes where local frustration
    makes isolated excitations energetically unfavorable.
    Superconductivity then emerges
    as the collective phase-locking
    of frustration-stabilized composites.

    Such an interpretation unifies mediator-driven
    and strongly correlated superconductors
    within a single structural framework,
    while remaining falsifiable
    through symmetry selection,
    stiffness scaling,
    and disorder response.

    \bigskip
    Future work should aim at
    deriving the fiber cost functional
    from explicit microscopic models,
    quantifying higher-order corrections to $r_{\mathcal F}$ scaling,
    and extending the framework
    to other unconventional superconductors,
    including iron-based compounds and nickelate variants.
